\section{Results, discussion, limitations, and conclusion}
\subsection{Classification and sparse compression}
Expression has the strongest observed original test performance on all four
metrics (Table~\ref{tab:test}). The dictionary compresses 2,000 input genes to
64 coordinates, a 31.25-fold dimensional reduction, while retaining high
accuracy. Its small validation Macro-F1 advantage (0.984580 versus 0.984365)
does not persist on test. This is not evidence of superiority or equivalence;
we do not retune using test outcomes.
\begin{table}[htbp]\centering
\caption{Original frozen evaluation on 1,275 test cells from the shared donor
pool. Higher F1/accuracy is better; lower log loss is better.}
\label{tab:test}
\begin{tabular}{lrrrr}\toprule
Representation & Macro-F1 & Accuracy & Balanced acc. & Log loss\\\midrule
Expression & 0.993913 & 0.994510 & 0.993687 & 0.019454\\
Dictionary & 0.980962 & 0.989020 & 0.977933 & 0.036989\\
scGPT & 0.962358 & 0.968627 & 0.955637 & 0.096535\\\bottomrule
\end{tabular}
\end{table}

The selected dictionary has no unused atoms in training or validation.
With activity defined as $|z_k|>10^{-6}$, validation codes are 17.02\% active
(10.89 atoms per cell); training codes are 15.30\% active. Validation
reconstruction MSE is 0.7826, or 0.7924 of the zero-reconstruction MSE.
This is a reconstruction diagnostic, not a classification result. Lower
penalties improve reconstruction in this grid without monotonically improving
classification (Appendix~\ref{app:tables}). The observed baseline strength
shows why comparison against expression itself is necessary.

\subsection{Donor separation and classifier scaling}
Table~\ref{tab:donor} summarizes the retrospective matched-panel study.
Expression and dictionary donor-mean Macro-F1 are close, while scGPT is lower.
Four donors and overlapping training folds do not support a broad significance
or equivalence claim. These values cannot be subtracted from the original
test scores to estimate a pure donor effect: grouping, gene policy, scaling,
regularization selection, and dictionary initialization also changed.
\begin{table}[htbp]\centering
\caption{Retrospective donor Macro-F1. Standardization is the primary protocol;
SD describes four donor scores after averaging dictionary seeds within donor.}
\label{tab:donor}
\begin{tabular}{lcc}\toprule
Representation & Standardized mean $\pm$ SD & Unscaled mean\\\midrule
Expression & $0.9246\pm0.0549$ & 0.9246\\
Dictionary & $0.9215\pm0.0488$ & 0.9220\\
scGPT & $0.8823\pm0.0473$ & 0.8440\\\bottomrule
\end{tabular}
\end{table}
The scale ablation changes scGPT Macro-F1 substantially, despite using the
same classifier family and separately tuning regularization. Dictionary
Macro-F1 changes little here; expression is already standardized upstream.
This supports treating feature scale as part of a fair linear-probing
protocol, without claiming that matched genes and scaling eliminate all
representation differences.

\subsection{Numerical accuracy and interpretation}
Controlled dictionary initialization yields within-donor Macro-F1 SDs of
0.0008--0.0030. Atom reproducibility is weaker: mean matched absolute cosine
is 0.606, top-20-gene Jaccard overlap 0.344, and sign-aligned training-code
correlation 0.572. Stable prediction therefore does not imply uniquely
recovered gene programs.

Across 48 donor-study dictionary fits, coding KKT residuals are at most
$1.1\times10^{-6}$ and the smallest observed Gram eigenvalue is about 0.00488.
These support accurate, conditionally unique encodings. Projected
dictionary-gradient norms remain 1.90--4.06 and early stopping occurs after
2--7 epochs: absence of warnings does not establish joint stationarity or
global optimality. All 640 donor-study classifier fits completed without
convergence warnings, and the 40 saved outer prediction sets were checked
against IDs, labels, and recomputed metrics.

The original fitted dictionary is directly inspectable. For example, atom 2
has large positive weights on CUZD1, CTRC, and PNLIPRP1 and elevated mean
activity in acinar cells; atom 5 emphasizes COL1A2, COL3A1, and COL1A1 and
is active in annotated activated stellate cells. These are observed
weight/activity associations (Appendix~\ref{app:figures}), not formal
enrichment or biological validation. Signs and indices are run-specific.
No causal pathway claim follows, particularly given moderate atom stability.

\subsection{Limitations and conclusion}
The original split shares donors, and the secondary study is retrospective,
uses only four donors from one study, and inherits earlier choices. Some
held-out donors contain only two or three acinar cells; rare-class exclusion
further restricts the estimand. Annotation quality is assumed. The small
custom scGPT model, missing panel genes, distinct input transforms, and
unaudited pretraining overlap limit attribution to architecture or
pretraining. The adaptive validation search and three initialization seeds
provide limited coverage. Formal pathway enrichment, external-dataset
confirmation, and comprehensive hyperparameter or seed sweeps were not done.

For this cohort and linear-probing protocol, the supplied scGPT representation
does not improve over processed expression or sparse dictionary codes.
Dictionary learning offers sparse compression and inspectable associations,
but not established biological identifiability. The strongest methodological
lesson is to distinguish coding accuracy, predictive stability, and
biological reproducibility, and to control both donor separation and
regularization geometry. A future confirmatory study requires an untouched
cohort, a protocol fixed in advance, and a pretraining-overlap audit.
\label{main-report-end}
\FloatBarrier
